\section{问题定义}\label{sec:problem}

给定一条新闻或社交媒体消息 $x$（仅文本），闭卷虚假新闻检测要求在\textbf{不检索任何外部证据}的条件下预测其真实性标签 $y\in\{0,1\}$（$1$ 表示虚假）。训练集 $\mathcal{D}=\{(x_i,y_i)\}$ 只含新闻级标签。除新闻级概率 $\hat y=P(y{=}1\mid x)$ 外，我们还希望模型给出：（i）新闻中各条声明存在问题的概率，用于定位；（ii）对 LLM 在每条声明上的分析是否可信的估计；（iii）在内容不足以判断时拒答的能力。

本文把 $x$ 视为一个由若干\emph{实例}组成的包（bag），沿用多示例学习（multiple instance learning, MIL）的标准假设\citep{dietterich1997solving,carbonneau2018mil}：记隐变量 $z_k\in\{0,1\}$ 表示第 $k$ 个实例是否存在问题，则
\begin{equation}
y=1 \iff \exists k,\ z_k=1 .\label{eq:mil}
\end{equation}
该假设蕴含一个不对称性：真新闻（$y{=}0$）的所有实例都满足 $z_k{=}0$，而假新闻（$y{=}1$）只保证至少一个实例有问题。\ref{sec:loss} 节利用这一不对称性，在没有任何声明级标注的情况下为门控构造监督信号。假设何时不成立见 \ref{sec:discussion} 节的讨论。

\section{方法}\label{sec:method}

\subsection{总体框架}

\method{}（\textbf{C}laim-\textbf{L}evel \textbf{A}rgument \textbf{R}eliability \textbf{A}djudication）由一个冻结的 LLM 与一个可训练的小模型裁决器组成（图~\ref{fig:framework}）。LLM 对每条新闻只调用一次，输出结构化的“声明分解 + 正反论证”；裁决器读取新闻原文与 LLM 输出，在声明级决定采信 LLM 的程度，并聚合为新闻级判定。LLM 的输出可离线缓存，训练与推理阶段裁决器不再访问 LLM。

\begin{figure}[t]
\centering
\resizebox{\linewidth}{!}{% Framework figure (TikZ); natural width ~172 mm, scaled to \linewidth where it is included.
% Orange = LLM-side quantities, blue = SLM-side, green = gates.
\definecolor{cLLM}{HTML}{EB6834}
\definecolor{cSLM}{HTML}{2A78D6}
\definecolor{cGate}{HTML}{1BAF7A}
\definecolor{cInk}{HTML}{0B0B0B}
\definecolor{cMuted}{HTML}{6B6A66}
\definecolor{cLine}{HTML}{C3C2B7}
\begin{tikzpicture}[
  x=1mm, y=1mm, font=\footnotesize, >={Stealth[length=1.6mm]},
  box/.style={draw=cLine, rounded corners=2pt, fill=white, align=center, inner sep=2.5pt, text=cInk},
  llm/.style={box, draw=cLLM!80, fill=cLLM!7},
  slm/.style={box, draw=cSLM!80, fill=cSLM!7},
  gate/.style={box, draw=cGate!90, fill=cGate!9},
  flow/.style={->, draw=cMuted, line width=0.5pt},
  sup/.style={->, draw=cMuted, dashed, line width=0.5pt},
  lab/.style={font=\scriptsize, text=cMuted, inner sep=1.5pt}]

% ---------------- column A: news and LLM
\node[box, minimum width=31mm] (x) at (0,0) {新闻 $x$\\[-1pt]{\scriptsize 句子 $s_1,\dots,s_n$}};
\node[llm, minimum width=31mm] (llm) at (0,-19) {冻结的 LLM\\[-1pt]{\scriptsize 离线生成，结构化解码}};
\node[llm, minimum width=31mm, align=left] (out) at (0,-40.5) {%
  {\scriptsize 声明 $c_k$（类型 $t_k$，句 $j_k$）}\\[-1pt]
  {\scriptsize 论证 $a_k^{+},a_k^{-}$，怀疑度 $v_k$}\\[-1pt]
  {\scriptsize 整体判断（$k=0$）}\\[-1pt]
  {\scriptsize 论证前评分 $\tilde v_0$}};
\draw[flow] (x) -- (llm);
\draw[flow] (llm) -- (out);

% ---------------- column B: evidence and gates
\node[slm, minimum width=45mm] (sx) at (47,0) {整体文本分支\\[-1pt]{\scriptsize $s_x=\mathrm{MLP}_x(\bm h_x)$}};
\node[slm, minimum width=45mm] (sslm) at (47,-19) {SLM 证据\\[-1pt]{\scriptsize $s_k^{\mathrm{slm}}=\mathrm{MLP}_s(\bm c_k)$}};
\node[llm, minimum width=45mm] (sllm) at (47,-31.5) {LLM 证据\\[-1pt]{\scriptsize $s_k^{\mathrm{llm}}=\mathrm{MLP}_a(\Delta\bm u_k)+\alpha\,\mathrm{logit}\,v_k$}};
\node[gate, minimum width=45mm] (g) at (47,-44) {判定门控\\[-1pt]{\scriptsize $g_k$：LLM 证据的权重}};
\node[gate, minimum width=45mm] (r) at (47,-56.5) {相关性门控\\[-1pt]{\scriptsize $r_k$：该实例是否属于这条新闻}};

% ---------------- columns C, D: combination
\node[box, minimum width=40mm] (pk) at (97,-37.75) {实例概率\\[-1pt]{\scriptsize $p_k=\sigma\big(b_0+r_k(s_k^{\mathrm{slm}}+g_k\,s_k^{\mathrm{llm}})\big)$}\\[-1pt]{\scriptsize $r_k\to0$：空实例，$p_k\to\sigma(b_0)$}};
\node[box, minimum width=41mm] (nor) at (143,-37.75) {noisy-OR 聚合\\[-1pt]{\scriptsize $P_{\mathrm{MIL}}=1-\prod_k(1-p_k)$}};
\node[box, minimum width=41mm] (y) at (143,0) {新闻级预测\\[-1pt]{\scriptsize $\hat y=\sigma\big(s_x+\tilde\gamma\,[\mathrm{logit}\,P_{\mathrm{MIL}}-\mathrm{logit}\,P_{\varnothing}]+b\big)$}};

% ---------------- flows
\draw[flow] (x) -- (sx);
\coordinate (bus) at (20,-40.5);
\draw[draw=cMuted, line width=0.5pt] (out.east) -- (bus);
\draw[flow] (bus) |- (sslm.west);
\draw[flow] (bus) |- (sllm.west);
\draw[flow] (bus) |- (g.west);
\draw[flow] (bus) |- (r.west);
\coordinate (join) at (73,-37.75);
\draw[draw=cMuted, line width=0.5pt] (sslm.east) -| (join);
\draw[draw=cMuted, line width=0.5pt] (sllm.east) -| (join);
\draw[draw=cMuted, line width=0.5pt] (g.east) -| (join);
\draw[draw=cMuted, line width=0.5pt] (r.east) -| (join);
\draw[flow] (join) -- (pk.west);
\draw[flow] (pk) -- (nor);
\draw[flow] (nor) -- (y);
\draw[flow] (sx) -- (y);

% ---------------- instance plate
\begin{scope}[on background layer]
  \node[draw=cLine, dashed, rounded corners=3pt, inner sep=3pt, fit=(sslm)(sllm)(g)(r)(pk)] (plate) {};
\end{scope}
\node[lab, anchor=south east] at ([xshift=-1mm,yshift=0.5mm]plate.south east) {对每个实例 $k=0,\dots,K$};

% ---------------- supervision (training only)
\node[lab] (lcls) at (143,11) {$\mathrm{BCE}(\hat y,\,y)$};
\node[lab] (lmil) at (143,-50) {$\mathrm{BCE}(P_{\mathrm{MIL}},\,y)$};
\node[lab, align=left] (lgate) at (97,-54) {$\mathcal L_{\mathrm{gate}}$：LLM 是否站对边\\[-1pt]$\mathcal L_{\mathrm{rel}}$：错配负例};
\draw[sup] (lcls) -- (y);
\draw[sup] (lmil) -- (nor);
\draw[sup] (lgate.west) -- ++(-6,0) |- ([yshift=-2mm]g.east);
\draw[sup] (lgate.west) -- ++(-6,0) |- ([yshift=-2mm]r.east);
\end{tikzpicture}
}
\caption{\method{} 总体框架。LLM 单次调用产生整体实例（$k{=}0$）与 $K$ 条声明，每个实例附带支持论证 $a^+$、质疑论证 $a^-$ 与怀疑度 $v$。裁决器对每个实例计算 SLM 证据、LLM 证据与可靠性门控，经 noisy-OR 聚合后与整体文本分支融合。$\Delta\bm u_k$ 表示两方论证表示的差与逐元素积，$\tilde\gamma=\softplus(\gamma)$。虚线箭头表示只在训练时使用的监督信号。}
\label{fig:framework}
\end{figure}

\subsection{单次调用的声明分解与双立场论证}\label{sec:llm}

\paragraph{输出结构}
先用确定性规则把 $x$ 切分为句子 $s_1,\dots,s_n$（$n\le 12$）并编号。LLM 读取编号后的文本，按固定顺序生成一个 JSON 对象：
\begin{enumerate}
\item \textbf{第一印象} $v^{\mathrm{fi}}$：在做任何分析之前，对“整条消息是虚假或误导性信息”的 1--5 级评分；
\item \textbf{声明列表} $\{c_k\}_{k=1}^{K}$（$K\le 6$）：每条声明包含自包含的声明文本、类型 $t_k\in\{$事实、数字、引述、观点$\}$、所在句子编号 $j_k$、认为其可信的最强理由 $a_k^+$、认为其可疑的最强理由 $a_k^-$，以及 1--5 级怀疑度；
\item \textbf{整体判断}：对整条消息的正反理由 $a_0^+,a_0^-$ 与 1--5 级怀疑度。
\end{enumerate}
整体判断在裁决器中作为第 $0$ 个实例，用于覆盖标题与正文不符、来源不明、情绪煽动等无法归到单条声明上的问题。完整提示词见附录~\ref{app:prompt}。

\paragraph{设计动机}
与多轮辩论\citep{liu2025ted,han2025d2d}相比，这一设计有三点考虑。其一，把论证单位从“发言轮次”换成“声明”，每个声明都有成对的正反论证，裁决器可以直接比较两方；LLM 调用次数固定为 1。其二，要求 LLM 先写两方论证、最后才给评分，并禁止在论证中写结论，以减轻先下结论再找理由的立场锁定。其三，提示词要求保留归属（“某人称……”不得改写为既成事实），因为一条真新闻可以如实转述一个不实说法，丢掉归属会直接破坏式~\eqref{eq:mil} 的假设。

\paragraph{结构化解码与软评分}
生成时用 JSON schema 约束解码\citep{kwon2023vllm,dong2024xgrammar}：声明数下限随句子数增加（$n\le2$ 时为 1，$n\le4$ 时为 2，否则为 3），句子编号限制在 $\{1,\dots,n\}$ 内。评分不直接使用生成的整数，而是取评分位置上五个等级 token 的概率 $\pi_\ell$ 并归一化，计算期望
\begin{equation}
v=\sum_{\ell=1}^{5}\frac{\ell-1}{4}\,\pi_\ell\in[0,1],
\end{equation}
同时保留归一化熵 $H(\pi)$ 作为 LLM 自身不确定性的度量。

\paragraph{论证顺序平衡}
LLM 的评分可能受最后写下的论证影响。我们按样本编号的奇偶确定性地交替 $a^+$ 与 $a^-$ 的书写顺序，使这一潜在偏置与标签无关。

\subsection{编码}\label{sec:enc}

所有方法遵循 ARG\citep{hu2024bad} 的设置：预训练编码器只微调最后一层，其余层冻结。\method{} 使用两份最后一层的可训练副本，$E_x$ 编码新闻原文，$E_a$ 编码 LLM 侧文本。记 $\operatorname{Pool}$ 为可学习的注意力池化，则
\begin{align}
\bm H_x &= E_x(x), & \bm h_x &= \operatorname{Pool}(\bm H_x),\\
\bm c_k &= \operatorname{Pool}\!\big(E_a(c_k)\big), & \bm u_k^{\pm} &= \operatorname{Pool}\!\big(E_a([c_k;a_k^{\pm}])\big),\\
\bm e_k &= \operatorname{LN}\!\big(\bm c_k+\operatorname{CrossAttn}(\bm c_k,\bm H_x)\big),
\end{align}
其中 $[c_k;a_k^{\pm}]$ 表示把声明与论证作为句对输入，$\bm e_k$ 是声明以自身为查询、对新闻 token 做交叉注意力后得到的上下文化表示。整体实例（$k{=}0$）的声明文本取一个固定短语。

由于冻结层的输出与训练轮次无关，我们对每条序列只做一次冻结层前向并缓存其隐状态，训练时仅计算最后一层。这与每步重新计算的差别仅在于冻结层没有 dropout 噪声；附录~\ref{app:live} 核对了两者的一致性。

\subsection{实例打分：SLM 证据与 LLM 证据}\label{sec:score}

对每个实例计算两路证据：
\begin{align}
s_k^{\mathrm{slm}} &= \operatorname{MLP}_s(\bm c_k),\\
s_k^{\mathrm{llm}} &= \operatorname{MLP}_a\!\big([\bm u_k^- - \bm u_k^+;\ \bm u_k^-\odot\bm u_k^+]\big)+\alpha\,\logit(v_k)+\mathbb{1}[k{=}0]\,\alpha_0\,\logit(v^{\mathrm{fi}}),
\end{align}
其中 $\alpha,\alpha_0$ 为可学习标量。$s^{\mathrm{llm}}$ 中的差分与逐元素积显式刻画两方论证的不对称，即“哪一方更容易被辩护”。

\paragraph{实例特定的分解式结构}
两路证据都只依赖该实例自身的输入：声明文本、它的两条论证与评分，而不依赖新闻整体表示 $\bm h_x$。这是有意为之。若实例打分可以看到整条新闻，noisy-OR 型 MIL 存在一个平凡解：每个实例都输出新闻级判断的副本。此时新闻级精度不受影响，但 $p_k$ 失去实例含义，定位无从谈起。新闻级的文本证据另有通道（\ref{sec:agg} 节的 $s_x$）。实验中我们把允许实例打分读取 $\bm e_k$ 与 $\bm h_x$ 的耦合变体作为消融对照。

\subsection{声明级可靠性门控}\label{sec:gate}

LLM 在闭卷条件下只能依靠参数记忆与常识。它对某些声明的分析有信息量，对另一些（例如需要实时事实的数字）则可能是幻觉。门控估计“LLM 在该实例上的分析是否可信”：
\begin{equation}
g_k=\sigma\!\Big(\operatorname{MLP}_g\big([\bm e_k;\ \bm h_x;\ \bm u_k^+;\ \bm u_k^-;\ \operatorname{emb}(t_k);\ \bm\phi_k]\big)\Big),
\end{equation}
其中 $\bm\phi_k=[v_k,\ |v_k-0.5|,\ \tilde v_k,\ |\tilde v_k-v_k|,\ H(\pi_k)]$ 是由评分导出的特征，$\tilde v_k$ 对整体实例取第一印象 $v^{\mathrm{fi}}$、对声明取 $v_k$。门控只缩放 LLM 证据，因此允许它读取上下文：
\begin{equation}
p_k=\sigma\!\big(s_k^{\mathrm{slm}}+g_k\,s_k^{\mathrm{llm}}\big).\label{eq:pk}
\end{equation}
$g_k\to0$ 时该实例退回到 SLM 自己的判断。

\subsection{Noisy-OR 聚合与整体分支}\label{sec:agg}

在实例相互独立的近似下，式~\eqref{eq:mil} 对应 noisy-OR 聚合\citep{maron1998framework}：
\begin{equation}
P_{\mathrm{MIL}}=1-\prod_{k=0}^{K}(1-p_k).
\end{equation}
MIL 实例只能把预测推向“假”，而新闻原文的整体表示在已有工作中贡献最大\citep{he2026factguard}，因此保留一个可以双向调节的整体文本分支：
\begin{equation}
\hat y=\sigma\!\big(s_x+\softplus(\gamma)\cdot\logit(P_{\mathrm{MIL}})+b\big),\qquad s_x=\operatorname{MLP}_x(\bm h_x),
\end{equation}
其中 $\gamma,b$ 为可学习标量。

\subsection{训练目标与门控的弱监督}\label{sec:loss}

总损失为
\begin{equation}
\mathcal L=\operatorname{BCE}(\hat y,y)+\lambda_{\mathrm{mil}}\operatorname{BCE}(P_{\mathrm{MIL}},y)+\lambda_{\mathrm{gate}}\,\mathcal L_{\mathrm{gate}} .
\end{equation}
当 $y{=}0$ 时，$\operatorname{BCE}(P_{\mathrm{MIL}},0)=-\sum_k\log(1-p_k)$，即对真新闻的每个实例施加了精确的负监督。

\paragraph{门控目标}
门控的目标是“LLM 在该实例上是否站对了边”。记 LLM 的表态分数为 $\bar v_k$（整体实例取 $(v_0+v^{\mathrm{fi}})/2$，声明取 $v_k$），$\bar v_k>0.5+\delta$ 视为判其有问题，$\bar v_k<0.5-\delta$ 视为判其无问题，介于两者之间视为未表态，不参与门控损失。于是
\begin{equation}
r_k=\tilde z_k\cdot\mathbb{1}[\bar v_k>0.5+\delta]+(1-\tilde z_k)\cdot\mathbb{1}[\bar v_k<0.5-\delta],
\end{equation}
其中 $\tilde z_k$ 是对隐状态 $z_k$ 的估计，分三种情况：
\begin{itemize}
\item \textbf{整体实例}：$\tilde z_0=y$，可直接观测；
\item \textbf{真新闻中的声明}：由 MIL 不对称性，$\tilde z_k=0$，即 LLM 判对当且仅当它没有怀疑该声明；
\item \textbf{假新闻中的声明}：$z_k$ 不可观测，用 noisy-OR 后验 $\tilde z_k=\sg\!\big[p_k/P_{\mathrm{MIL}}\big]$ 作为软责任（$\sg$ 表示停止梯度），在 $T_w$ 轮 warm-up 之后启用，并以权重 $\eta$ 计入。
\end{itemize}
$\mathcal L_{\mathrm{gate}}$ 是 $\operatorname{BCE}(g_k,r_k)$ 在已表态实例上的加权平均。与 ARG、FactGuard 的实例级可用性监督（整条理由只有一个“LLM 判断是否等于标签”的比特）相比，上述目标把监督下沉到了声明。

算法~\ref{alg:train} 给出完整的训练流程。

\begin{algorithm}[t]
\caption{\method{} 的训练}\label{alg:train}
\KwIn{训练集 $\mathcal D$；冻结的 LLM；超参数 $\lambda_{\mathrm{mil}},\lambda_{\mathrm{gate}},\eta,\delta,T_w$}
\KwOut{裁决器参数 $\theta$}
\ForEach{$x\in\mathcal D_{\mathrm{train}}\cup\mathcal D_{\mathrm{val}}\cup\mathcal D_{\mathrm{test}}$}{
  $A(x)\leftarrow\operatorname{LLM}(x)$\tcp*{一次调用，结果缓存}
}
缓存所有序列的冻结层隐状态\;
\For{$\mathrm{epoch}=0,1,\dots$}{
  \ForEach{小批量 $(x,y,A)$}{
    计算 $\bm h_x,\bm c_k,\bm e_k,\bm u_k^\pm$，以及 $s_k^{\mathrm{slm}},s_k^{\mathrm{llm}},g_k,p_k$\;
    $P_{\mathrm{MIL}}\leftarrow1-\prod_k(1-p_k)$；$\hat y\leftarrow\sigma(s_x+\softplus(\gamma)\logit(P_{\mathrm{MIL}})+b)$\;
    按 \ref{sec:loss} 节构造门控目标 $r_k$ 与权重；若 $\mathrm{epoch}<T_w$，假新闻声明的权重置 0\;
    $\mathcal L\leftarrow\operatorname{BCE}(\hat y,y)+\lambda_{\mathrm{mil}}\operatorname{BCE}(P_{\mathrm{MIL}},y)+\lambda_{\mathrm{gate}}\mathcal L_{\mathrm{gate}}$；更新 $\theta$\;
  }
  在验证集上计算 macro-F1，按早停规则保留最优参数\;
}
\end{algorithm}

\subsection{\student{}：声明到句子对齐的蒸馏}\label{sec:student}

为了在推理时完全不调用 LLM，我们训练一个只读新闻的学生模型。学生以新闻的\emph{句子}为 MIL 实例（外加一个整体实例），沿用教师的 SLM 分支：$p_j^S=\sigma(\operatorname{MLP}_s(\bm c_j))$，聚合与融合方式同 \ref{sec:agg} 节。LLM 为每条声明给出的句子编号 $j_k$ 把教师的声明级概率对齐到句子：
\begin{equation}
t_j=\max_{k:\,j_k=j}p_k^{T},
\end{equation}
没有声明与之对齐的句子不参与实例级蒸馏。学生的损失为
\begin{equation}
\mathcal L_S=\operatorname{BCE}(\hat y^S,y)+\lambda_{\mathrm{mil}}\operatorname{BCE}(P^S_{\mathrm{MIL}},y)+\kappa_{o}\operatorname{BCE}(\hat y^S,\hat y^T)+\kappa_{i}\,\overline{\operatorname{BCE}}(p^S_j,t_j)+\kappa_{f}\operatorname{MSE}(\bm h^S_x,\bm h^T_x),
\end{equation}
参数由教师的对应模块初始化。ARG-D 与 FactGuard-D 只蒸馏池化后的特征向量，学生没有实例级输出；\student{} 的学生保留了句级定位能力。

\subsection{选择性预测}\label{sec:selective}

闭卷条件下部分新闻无法从内容判定。除常用的 MaxProb 置信度 $|\hat y-0.5|$ 外，\method{} 的中间量提供了额外的“不可判”信号。我们在验证集上拟合一个逻辑回归选择器，预测裁决器自身的判定是否正确，输入特征为 $|\hat y-0.5|$、两分支分歧 $|P_{\mathrm{MIL}}-\sigma(s_x)|$、门控均值、$\max_k p_k$、$|P_{\mathrm{MIL}}-0.5|$ 与实例数的对数。选择器不接触测试标签。

\subsection{开销}\label{sec:cost}

每条新闻需要一次 LLM 调用，输出长度与声明数成正比。裁决器在 BERT-base 规模的编码器上增加一份最后一层副本与若干 MLP。每条新闻需要编码 $1+3(K{+}1)$ 条短序列（新闻、每个实例的声明及两个声明--论证对），其冻结层前向只做一次。实测开销见 \ref{sec:exp-cost} 节。
